% Options for packages loaded elsewhere
\PassOptionsToPackage{unicode}{hyperref}
\PassOptionsToPackage{hyphens}{url}
\documentclass[
]{article}
\usepackage{xcolor}
\usepackage{amsmath,amssymb}
\setcounter{secnumdepth}{5}
\usepackage{iftex}
\ifPDFTeX
  \usepackage[T1]{fontenc}
  \usepackage[utf8]{inputenc}
  \usepackage{textcomp} % provide euro and other symbols
\else % if luatex or xetex
  \usepackage{unicode-math} % this also loads fontspec
  \defaultfontfeatures{Scale=MatchLowercase}
  \defaultfontfeatures[\rmfamily]{Ligatures=TeX,Scale=1}
\fi
\usepackage{lmodern}
\ifPDFTeX\else
  % xetex/luatex font selection
\fi
% Use upquote if available, for straight quotes in verbatim environments
\IfFileExists{upquote.sty}{\usepackage{upquote}}{}
\IfFileExists{microtype.sty}{% use microtype if available
  \usepackage[]{microtype}
  \UseMicrotypeSet[protrusion]{basicmath} % disable protrusion for tt fonts
}{}
\makeatletter
\@ifundefined{KOMAClassName}{% if non-KOMA class
  \IfFileExists{parskip.sty}{%
    \usepackage{parskip}
  }{% else
    \setlength{\parindent}{0pt}
    \setlength{\parskip}{6pt plus 2pt minus 1pt}}
}{% if KOMA class
  \KOMAoptions{parskip=half}}
\makeatother
\usepackage{graphicx}
\makeatletter
\newsavebox\pandoc@box
\newcommand*\pandocbounded[1]{% scales image to fit in text height/width
  \sbox\pandoc@box{#1}%
  \Gscale@div\@tempa{\textheight}{\dimexpr\ht\pandoc@box+\dp\pandoc@box\relax}%
  \Gscale@div\@tempb{\linewidth}{\wd\pandoc@box}%
  \ifdim\@tempb\p@<\@tempa\p@\let\@tempa\@tempb\fi% select the smaller of both
  \ifdim\@tempa\p@<\p@\scalebox{\@tempa}{\usebox\pandoc@box}%
  \else\usebox{\pandoc@box}%
  \fi%
}
% Set default figure placement to htbp
\def\fps@figure{htbp}
\makeatother
\setlength{\emergencystretch}{3em} % prevent overfull lines
\providecommand{\tightlist}{%
  \setlength{\itemsep}{0pt}\setlength{\parskip}{0pt}}
\input{format}
\asignatura{Tipología y ciclo de la vida de los datos}
\actividad{PRACT1}
\titulo{Práctica 2}
\subtitulo{¿Cómo realizar la limpieza y análisis de los datos?}
\autor{\begin{tabular}[t]{@{}l@{}} Agustín Barbatelli Balboa\\ Adrián Cordones Martínez \end{tabular}}
\usepackage{bookmark}
\IfFileExists{xurl.sty}{\usepackage{xurl}}{} % add URL line breaks if available
\urlstyle{same}
\hypersetup{
  hidelinks,
  pdfcreator={LaTeX via pandoc}}

\author{}
\date{\vspace{-2.5em}}

\begin{document}

\sinportada

\criterios{Metadatos}{
  \item \textbf{Integrantes del grupo:}
    \begin{itemize}
      \item Agustín Barbatelli Balboa
      \item Adrián Cordones Martínez
    \end{itemize}
    
  \item \textbf{Enlace al repositorio con el código:}\\
    \url{https://github.com/adriancordones/steam-apps-analysis}
    
  \item \textbf{Enlace al vídeo de explicación:}\\
    \url{https://...}
}

\section{Descripción}\label{descripciuxf3n}

Este dataset recoge metadatos de aplicaciones publicadas en la \textit{Steam Store}, extraídos mediante web scraping de las páginas individuales de producto. Cada registro corresponde a un juego o aplicación e incluye información identificativa (título, desarrollador, editor), condiciones comerciales (precio, descuento), recepción por parte de los usuarios (número y resumen general de reseñas), características de contenido (DLCs, logros, idiomas, banda sonora) y clasificación temática (géneros y etiquetas de la comunidad). Los datos reflejan el estado de la tienda en el momento de la extracción (abril de 2026) y cubren los títulos más populares y recientes según los resultados de búsqueda de Steam para España y con la página en español.

\section{Integración y selección}\label{integraciuxf3n-y-selecciuxf3n}

...

\section{Limpieza de los datos}\label{limpieza-de-los-datos}

...

\section{Análisis de los datos}\label{anuxe1lisis-de-los-datos}

...

\section{Representación de los
resultados}\label{representaciuxf3n-de-los-resultados}

...

\section{Resolución del problema}\label{resoluciuxf3n-del-problema}

...

\section{Código}\label{cuxf3digo}

...

\section{Vídeo}\label{vuxeddeo}

...

\begin{thebibliography}{9}\label{ref:Ref}
  \bibitem{R1} De Luisa, A., Hartman, J., Nabergoj, D., Pahor, S., Rus, M., Stevanoski, B., Demšar, J. y Štrumbelj, E. (2021). Predicting the Popularity of Games on Steam. arXiv:2110.02896. URL: \url{https://arxiv.org/abs/2110.02896}
  \bibitem{R2} Teja, A. S., Hanafi, M. L. I. y Qomariyah, N. N. (2023). Predicting Steam Games Rating with Regression. E3S Web of Conferences, ICOBAR 2023. URL: \url{https://www.e3s-conferences.org/articles/e3sconf/pdf/2023/25/e3sconf_icobar2023_02001.pdf}
  \bibitem{R3} Creative Commons (2023) Licencia Creative Commons Atribución-NoComercial-CompartirIgual 4.0 Internacional. [En línea]. URL: \url{https://creativecommons.org/licenses/by-nc-sa/4.0/deed.es} [Consultado: abril de 2026].
  \bibitem{R4} Subirats Maté, L. y Calvo González, M. (2019) Web scraping. Barcelona: Editorial UOC.
  \bibitem{R5} Lawson, R. (2015) Web Scraping with Python: Successfully scrape data from any website with the power of Python. Birmingham: Packt Publishing. ISBN: 9781782164364.
  \bibitem{R6} Scrapy developers (2026) Scrapy documentation (versión 2.14.2). [En línea]. URL: \url{https://docs.scrapy.org/en/latest/} [Consultado: abril de 2026].
  \bibitem{R7} Valve Corporation (2026) Steam Store. [En línea]. URL: \url{https://store.steampowered.com} [Consultado: abril de 2026].
\end{thebibliography}

\end{document}
